\documentclass[letterpaper]{article} % DO NOT CHANGE THIS
\usepackage[submission]{aaai2027}  % DO NOT CHANGE THIS
\usepackage[hyphens]{url}  % DO NOT CHANGE THIS
\usepackage{graphicx} % DO NOT CHANGE THIS
\urlstyle{rm} % DO NOT CHANGE THIS
\def\UrlFont{\rm}  % DO NOT CHANGE THIS
\usepackage{natbib}  % DO NOT CHANGE THIS
\usepackage{caption} % DO NOT CHANGE THIS
\captionsetup{font=small,skip=4pt}
\frenchspacing  % DO NOT CHANGE THIS
\usepackage{booktabs}
\usepackage{amsmath,amssymb}
\usepackage{tikz}
\usetikzlibrary{shapes.geometric,arrows.meta,positioning}
\pdfinfo{
/TemplateVersion (2027.1)
}
\setcounter{secnumdepth}{0}

\title{Opportunity Certification: Deciding Whether a Learned Component\\ Is Warranted, Before Training One}
\author{Anonymous Submission}
\affiliations{}

\newcommand{\x}{$\times$}
\newcommand{\orc}{$^{\dagger}$} % marks a clairvoyant (non-deployable) ingredient
\newcommand{\rung}[1]{\textsc{#1}}
% Use the external figure when it is present, otherwise a visible placeholder, so the
% document still compiles on a checkout that lacks the exported artwork.
\newcommand{\figorbox}[2]{%
  \IfFileExists{#1.pdf}{\includegraphics[#2]{#1.pdf}}{%
  \IfFileExists{#1.png}{\includegraphics[#2]{#1.png}}{%
  \fbox{\parbox[c][2.4cm][c]{0.9\columnwidth}{\centering\ttfamily\small [#1 not found]}}}}}

\begin{document}
\maketitle

\begin{abstract}
\noindent In machine learning for systems, a gap between a deployed heuristic and an oracle with future knowledge is often interpreted as an opportunity for a learned model to fill. This work tests that inference in advance of any model training, using cache prefetching as the demonstration domain. Much of a raw oracle advantage is due to weak baselines, unmetered bandwidth, and objects no deployable predictor could foresee, so the first question is how much of the gap exists in principle. A measurement procedure, the Instrument, eliminates all three: it compares an adversarially tuned baseline against a future-seeing policy, held to the same bandwidth and restricted to objects a past-trained predictor could have named. Across twelve production traces the apparent opportunity shrinks by $14$--$87\%$, surviving on three at $12.7$--$19.7$ hit-rate points. A surviving gap need not be a capturable one. The Necessity Ladder therefore grants the oracle one perfect capability at a time: if perfection in a capability cannot close the gap, no learner improving only that capability can. Perfect timing recovers $91$--$96\%$ of the corridor and clairvoyant coverage clears the bar, so the corridor is reachable in principle; every deployable ingredient then falls short. Perfect arbitration adds $\pm 0.01$ points, fetching at emission with a flawless forecaster recovers at most $6\%$, and a learned timing model lands $26$--$40$ points below the baseline despite ranking reuse accurately, a better forecaster (rank correlation $0.60$ to $0.68$) capturing nothing further. Ranking is not the binding constraint: holding coverage at the oracle's and granting perfect timing outright, a policy must still convert $58\%$ of the bytes it spends into hits to match the baseline, where the best causal model converts $6\%$. Narrowing does not rescue it, the precision so bought costing more coverage than it returns, so a signal can be rank-learnable without being budget-feasible; a policy built against this bottleneck also fails, its added signal proving redundant with the predictor's own confidence. Applied unchanged, the procedure returns no corridor for language-model key-value cache warming, where a tuned lookup table already sits on the physically reachable ceiling, and build for a learned index on one of five key sets, withholding the verdict wherever a tuned classical structure already saturates that ceiling. Two trained policies bound the false-negative rate. Four settings, four pre-registered verdicts, none requiring a trained model to reach.
\end{abstract}


\section{Introduction}
Often in machine learning for systems, a response to the errors made by a deployed heuristic is to measure the gap between the heuristic and an oracle and to attempt to train a model to capture the former as a surrogate for the latter. This inference rests, however, on two silent premises: that the gap was measured honestly, and that the information used by the oracle is information a deployable model could hold. This paper argues that both premises are violated often enough that the inference itself should be tested, and tested \emph{before} a model is trained. The argument is developed in cache prefetching, a domain in which both failures occur and can be measured, and the resulting procedure is then applied unchanged to two further domains.

The prevailing yardsticks invite mismeasurement. The gap to Belady's optimum \citep{belady} ignores prefetching, and its prefetch-aware repair \citep{demandmin} shows that the optimum itself moves once prefetching is admitted. Oracle ceilings spend unbounded bandwidth against throttled or under-tuned baselines, and clairvoyant ceilings credit \emph{compulsory} misses, first references to never-seen objects that no history-based policy can prefetch. In preliminary measurement for this work, a $+42$-point synthetic headroom collapsed to $-0.02$ once the baseline was tuned across a predictor family and the oracle was charged the baseline's exact prefetch bandwidth. That collapse fixed the two commitments everything below enforces: a strongest-available baseline, and every ceiling held to the baseline's byte budget, under pre-registered thresholds and machine-checked invariants.

The proposal, \textbf{Opportunity Certification}, submits an apparent oracle gap to two stages (Figure~\ref{fig:overview}). \emph{Stage 1, the Instrument,} certifies how much of the gap is real, measuring a \textbf{learnable corridor}: the distance between the best decoupled predictor over a swept family and a ceiling restricted to objects a history-based predictor could name, both at one prefetch byte budget. \emph{Stage 2, the Necessity Ladder,} certifies whether a real corridor is reachable, pricing each capability a learner could contribute with a separate oracle. The output is a certificate: \emph{no corridor}, a \emph{trap}, or \emph{build}.

On twelve production traces the Instrument deflates gross corridors by $14$--$87\%$; three survive at $+12.7$ to $+19.7$ hit-rate points, with bootstrap intervals clear of a pre-registered bar. Every deployable ingredient then falls short at a single measured price, a \textbf{break-even precision} of $0.582$ against $0.062$ achieved, so a reuse signal can be \emph{rank-learnable} without being \emph{budget-feasible}; a policy built specifically against that bottleneck also fails, its added signal proving redundant with the predictor's own confidence.

Two further instantiations show the procedure discriminates rather than merely deflates: \emph{no corridor} for key-value cache warming in LLM serving, where a tuned association table already sits on the physically reachable ceiling, and \emph{build} for a learned index on one of five key sets, the verdict withheld on three others where a tuned classical structure already attains that ceiling. Reachability separates a build from a trap; baseline saturation separates a build from an illusion.

\paragraph{Contributions.} (i) Opportunity Certification, a two-stage, domain-general procedure with pre-registered thresholds and machine-checkable invariants; (ii) the Instrument and a deflationary measurement across 12 traces, with a taxonomy of four inflation mechanisms and the finding that eviction and prefetching are substitutes; (iii) the Necessity Ladder and its threefold negative on the surviving corridors; (iv) a break-even precision that prices the shortfall independently of any particular timing model, separating rank-learnability from budget feasibility; (v) a cross-domain instantiation returning four distinct verdicts, with a captured corridor as a positive control and a trained-model competence control establishing that the substrate admits learned wins.

\begin{figure}[t]\centering
\figorbox{figures/overview}{width=\columnwidth}
\caption{\textbf{Opportunity Certification.} Stage 1 (the Instrument) certifies how much of an apparent oracle gap is real; Stage 2 (the Necessity Ladder) certifies whether a deployable policy can reach it. The certificate reads no corridor, trap, or build.}
\label{fig:overview}
\end{figure}

\section{Related Work}
\textbf{Optimal replacement and prefetch-aware repairs.} Belady's MIN \citep{belady} is optimal for demand caching without prefetching; \citet{demandmin} showed that under prefetching the optimum becomes bandwidth-relative. This work takes that insight to variable-size object caches and adds the reachability split: it bounds what may be done with a history-based policy, not what an omniscient one may. \citet{cao} considered integrated prefetching and caching under clairvoyance, uniform sizes, and a latency model, assumptions that fail for CDN and key-value objects, so ceilings here are measured per trace.

\textbf{Learned caching.} LRB \citep{lrb} imitates a relaxed Belady oracle for eviction; Cold-RL \citep{coldrl}, DEAP \citep{deap}, DeePref \citep{deepref}, and LightCacheRL \citep{lightcache} learn eviction or prefetching; Pythia \citep{pythia} tunes hardware prefetching with online RL; \citet{jointhw} share representations across replacement and prefetching. \textbf{Baleen} \citep{baleen} is the nearest neighbor and its reported gains are not disputed, because the measured quantity differs on four axes: it optimises disk-head time under a flash write budget rather than hit ratio at a tuned baseline's byte rate; its dominant lever is admission, which commits no speculative bytes and so faces no break-even precision of this kind; its OPT-derived bound is not restricted to a predictor's vocabulary (M4); and it is compared against production configurations rather than a family-swept baseline (M1). A gain there and a trap here can both be true. To current knowledge no prior work constructs a pre-registered, bandwidth-matched, learnability-split headroom measurement for object caches, nor gates a learned component on an oracle decomposition of necessity.

\textbf{LLM serving caches and learned indexes.} Mooncake \citep{mooncake} released the production KV-cache traces used below; PagedAttention \citep{vllm} and CacheGen \citep{cachegen} established and compressed prefix-shared KV caching. That literature reports achieved hit rates; the cross-domain section measures headroom. Learned indexes \citep{kraska,sosd} supply the positive foil. \citet{sober} showed claimed reasoning gains collapse under variance-aware protocols; this paper is the systems analog, adding a capturability test ahead of any training. S3-FIFO \citep{s3fifo} is the fixed evictor throughout, so every corridor is attributable to prefetch scheduling.


\section{Stage 1: The Instrument}
The Instrument is one deterministic replay in which four arms share a single prefetch byte budget, and every degree of freedom that could inflate the corridor is fixed before any real-trace run.

\textbf{Setup.} Each run replays a two-million-request prefix of a production trace through the same cache: capacity 1\% of the trace footprint, a 5\% warmup excluded from scoring, request-weighted hit ratio as the score, S3-FIFO as the evictor in every arm. Predictors are decoupled from the future: they train on the first half of the trace and are frozen. A prefetched object later requested counts as a hit, one evicted before use counts as waste, and every prefetched byte counts as traffic; the accounting is unit-tested.

\textbf{The four arms (supplementary document, \S F).} \emph{BASE} is S3-FIFO with no prefetch; its origin traffic is the denominator. \emph{TUNED BASELINE} is the best non-learned system constructible against the corridor's own interest: the highest hit ratio over $\{$Markov-1, Markov-2, Markov-3, LSTM$\}$ across all confidence thresholds and fan-outs, capped at $1.15\times$ BASE traffic. \emph{GROSS CEILING} is a clairvoyant prefetcher metered by a token bucket to the baseline's byte rate, reported only as a diagnostic because it fetches objects no real predictor could know. \emph{WARM CEILING}, the primary arm, is the same clairvoyant restricted to objects seen during training, at the same budget; whatever the gross oracle wastes on unreachable objects, the warm ceiling re-spends on reachable ones. It is the honest upper bound for a history-based scheduler: perfect timing, realistic knowledge, equal cost.

\textbf{Why this interval.} It is the height range above the strongest system one can perform without new ideas and below what any history-based policy could in principle achieve at that budget: the low tempting shortcut, learnable $\approx$ gross corridor minus cold slice, is wrong in both directions: 27\% of the oracle's useful prefetches on \texttt{Wikipedia} were cold, and redirecting that budget to be warm helps $+5.4$, so subtracting misstates the corridor by as many as 5.5 points, more than enough to flip one vote. So must the reachability price be paid within the replay, and the invariant that requires exactly zero cold prefetch hits within the warm arms of the baseline and the warm arm of the council enforces it, having already subsumed a friendlier ``cold'' definition that happily (and wrongly) awarded the baseline 442{,}188 impossible cold prefetch hits.

\textbf{Pre-registration.} The verdict bar, learnable corridor $\ge 8$ points \emph{and} warm-wise Pareto dominance at the warm ceiling, was written down on paper before any real-trace run. The baselines come in pre-registered stages (coarse Markov, then thresholds plus Markov-3, then the LSTM), and to pass the live verdict, one must pass all of them; and of course passing a stronger baseline can only raise the bar, so passing or not is meaningful and finally-decisive.


\begin{table*}[t]
\centering\small
\setlength{\tabcolsep}{4.5pt}
\begin{tabular}{llccccp{5.2cm}}
\toprule
trace & family & base & tuned baseline & gross corr. & verdict & reason for verdict \\
\midrule
\texttt{wiki\_2019t} & CDN & 0.152 & 0.5531 @1.09\x & $+26.49$ & retained & (none)  \\
\texttt{meta\_rprn} & CDN & 0.390 & 0.6738 @1.03\x & $+32.57$ & removed &  cold-dominated \\
\texttt{meta\_reag} & CDN & 0.538 & 0.7547 @1.11\x & $+24.51^{\ddagger}$ & removed &  cold-dominated \\
\texttt{cluster50} & KV & 0.416 & 0.7080 @1.02\x & $+20.37$ & retained & (none)  \\
\texttt{cluster53} & KV & 0.595 & 0.6101 @1.10\x & $+25.72^{\dagger}$ & retained & (none)  \\
\texttt{cluster26} & KV & 0.786 & 0.8907 @1.09\x & $+10.27$ & \textbf{removed} & traffic-bought (ceiling $1.16\times$ vs bar $1.09\times$) \\
\texttt{cluster10} & KV & 0.500 & 0.7368 @1.00\x & $-1.18$ & removed & degenerate predictability (baseline $>$ ceiling) \\
\texttt{msr\_proj\_0} & block & 0.552 & 0.7568 @1.13\x & $+4.16$ & removed & baseline saturation (robust to any cap) \\
\texttt{msr\_hm\_0} & block & 0.657 & 0.8380 @1.13\x & $+10.08$ & removed &  cold $+$ saturation \\
\texttt{msr\_web\_2} & block & 0.010 & 0.8284 @1.00\x & $-3.81$ & removed & baseline saturation (prec 0.996) \\
\texttt{w105} & block & 0.269 & 0.8494 @1.05\x & $-0.45$ & removed & baseline saturation \\
\texttt{w87} & block & 0.095 & 0.5383 @1.12\x & $+32.67$ & removed$^{*}$ & Pareto fail by $0.01\times$ (boundary) \\
\bottomrule
\end{tabular}
\caption{Gate A across 12 traces. $^{*}$\texttt{w87} fails Pareto dominance by $0.01\times$, inside display resolution; the pre-registered rule removes it, and no claim depends on it. $^{\dagger}$\texttt{cluster53} is quoted against its strongest bar, the LSTM; Table~\ref{tab:learnable} decomposes it against its Markov bar, giving $+24.46$, and the verdict is identical either way. $^{\ddagger}$\texttt{meta\_reag} differs by $0.04$ between tables, a re-run at identical settings; the smaller value is used wherever a corridor is claimed.}
\label{tab:gatea}
\end{table*}



\section{Results: Measurement}
\textbf{Which gross corridors survive.} Table~\ref{tab:gatea} shows Gate~A. Three observations are salient. The tuned baseline is all-important: \texttt{cluster26} falls below the threshold once the baseline escalates. The LSTM never recovers a corridor, so the survivors are not artifacts of an ineffectual predictor. And survival is not a family trait: block traces give both $+10.08$ and $-3.81$, key-value traces $+25.72$ and $-1.18$.

\textbf{14--87\% of the gross is phantom.} Table~\ref{tab:learnable} unpacks the six gross-live traces. The three with most spectacular gross corridors lose 51--87\% to the cold slice, since the high object churn exhausts the clairvoyant budget on compulsory misses; the two low-churn KV traces lose only 14--20\%. Wikipedia is the apprising example: phantom to the tune of 52\%, but still a fair amount of exploitable warm reuse sufficient for the honest corridor to clear the bar, however only after the warm ceiling re-spends the cold budget; a na\"ive subtraction would yield $+7.27$ and wrongly dismiss it. Corridors are paired mean differences over the same post-warmup requests, resampled as 1{,}000 contiguous blocks with 10{,}000 resamples (block, not i.i.d., because hits are autocorrelated through cache state). All three lower bounds clearly exceed 8. (\texttt{cluster53}'s decomposition uses its Markov baseline; against its higher LSTM baseline the corridor is $\ge +18.4$, therefore the verdict remains.)

\begin{table}[t]
\centering\footnotesize
\setlength{\tabcolsep}{2pt}
\begin{tabular}{lcccc}
\toprule
trace & gross & cold & phantom & \textbf{learnable} (CI) \\
\midrule
\texttt{cluster53} & $+24.46$ & 4.62 & 20\% & $\mathbf{+19.67}$ $[19.4,19.9]$ \\
\texttt{cluster50} & $+20.37$ & 2.69 & 14\% & $\mathbf{+17.56}$ $[15.8,19.4]$ \\
\texttt{wiki\_2019t} & $+26.49$ & 19.22 & 52\% & $\mathbf{+12.72}$ $[10.9,14.5]$ \\
\texttt{msr\_hm\_0} & $+10.08$ & 5.82 & 51\% & $+4.89$ (removed) \\
\texttt{meta\_rprn} & $+32.57$ & 27.90 & 87\% & $+4.32$ (removed) \\
\texttt{meta\_reag} & $+24.55$ & 20.71 & 85\% & $+3.63$ (removed) \\
\bottomrule
\end{tabular}
\caption{Warm/cold decomposition on the gross-live traces; every run verifies that bar-cold and warm-cold are both zero. Pre-registered rule: reliably live only if the CI lower bound is at least 8.}
\label{tab:learnable}
\end{table}


\textbf{The verdicts do not depend on the exact bar.} Pre-registration is a process claim a reader cannot audit, so it is supplemented with a robustness result that can be: sweeping the threshold from 4 to 12 points, applied as the rule applies it, to the bootstrap lower bound.

\vspace{2pt}\noindent{\small
\begin{tabular}{lcccccc}
\toprule
trace & CI lo & 4 & 6 & \textbf{8} & 10 & 12 \\
\midrule
\texttt{cluster53} & 19.4 & live & live & \textbf{live} & live & live \\
\texttt{cluster50} & 15.8 & live & live & \textbf{live} & live & live \\
\texttt{wiki\_2019t} & 10.9 & live & live & \textbf{live} & live & dead \\
\texttt{msr\_hm\_0} & 4.89$^{*}$ & live & dead & \textbf{dead} & dead & dead \\
\texttt{meta\_rprn} & 4.32$^{*}$ & live & dead & \textbf{dead} & dead & dead \\
\texttt{meta\_reag} & 3.63$^{*}$ & dead & dead & \textbf{dead} & dead & dead \\
\bottomrule
\end{tabular}}\vspace{2pt}

\noindent($^{*}$point estimates, conservative here since a bootstrap bound could only be smaller.) Every survivor clears 10 and every removal fails at 6. The other thresholds have wider margins still: \rung{r2}'s arbitration gap measured $\pm0.01$ against a 2-point requirement and so survives any bar above $0.1$, while in the index domain \texttt{meta\_rprn} would flip at 51\%, which is why it is reported as marginal. Whether the thresholds were fixed in advance therefore decides no verdict except that one.

\textbf{Four inflation mechanisms, caught on the record.} The deflation decomposes into four mechanisms, each caught by a named component and each first exposed on this work's own results rather than on the literature's. \textbf{M1}, an under-tuned baseline (synthetic $+42\to-0.02$; \texttt{cluster26} $+12.5$ to below bar), is removed by the family sweep over $\{$Markov-1/2/3, LSTM$\}\times\tau\times k$. \textbf{M2}, bandwidth mismatch (a synthetic oracle at $1.27\times$ baseline traffic), is removed by the token bucket at the baseline's byte rate. \textbf{M3}, a traffic-bought corridor (\texttt{cluster26}: ceiling $1.16\times$ against bar $1.09\times$), is removed by the Pareto-dominance check. \textbf{M4}, cold conflation (\texttt{meta\_rprn} $+32.6$ gross to $+4.3$ learnable), is removed by restricting the ceiling to in-vocabulary reuse; the four are tabulated in the supplementary document, \S D. Stage~2 adds five further self-caught errors, two inside the oracle arms themselves, and the sharpest instance, which both removed and restored verdicts under one rule, is reported with the learned-index results below. A protocol's credibility rests on what it catches, including from its own authors; a test that never overturns anything provides no evidence.


\textbf{What separates live from dead.} Where no corridor survives, the tuned baseline already reaches the warm ceiling (\texttt{meta\_reag} 0.754 versus 0.791), so the residue is compulsory-cold and unreachable by construction; where one survives, in-vocabulary reuse exists that the tuned predictor fails to schedule (Wikipedia 0.553 versus 0.680), which is precisely the quantity a learned scheduler would target.

\textbf{Eviction and prefetching are substitutes.} A $2\times2$ of $\{$S3-FIFO, Belady$\}\times\{$no prefetch, clairvoyant at matched bandwidth$\}$ tests their interaction. All six traces are negative: $-1.34$ to $-7.94$. Better prefetching reduces what better eviction can add on every trace measured, so decoupled scheduling is the indicated design choice, and this also justifies the decision to hold S3-FIFO fixed throughout Stage~2.

\textbf{Robustness.} The corridor shrinks monotonically as the cache grows, from $+15.2$ at a 0.1\% cache to $+12.7$ at 1\% and $+5.1$ at 10\% on Wikipedia, so survival is a property of the operating point: two further traces clear the bar at 0.1\% and none at 10\%. Byte-weighting deflates it again, sometimes below zero where the request-weighted figure is large (\texttt{meta\_reag} reads $+26.7$ by request and $+4.9$ by byte at 0.1\%). Since the surviving corridor therefore belongs to the request-weighted objective that latency-serving tiers are operated against, not to egress cost, an operator optimising bytes should read these traces as having no corridor at all. Stage~2's verdicts are unaffected, being negative under the more favourable weighting.


\section{Stage 2: The Necessity Ladder}
Three certified corridors invite the expected next step, a learned scheduler. The Ladder turns that impulse into a question sequence: each rung is a pre-registered, replay-only experiment that tests one competence with an oracle and sets its price prior to any model training (Figure~\ref{fig:ladder}, Table~\ref{tab:decomp}). Each arm holds S3-FIFO fixed, follows the baseline's byte budget via the identical token bucket, and gives paired block-bootstrap intervals; the oracle arms bound what is achievable and are never reported as deployable results. Each rung grants strictly less foresight than the previous one, so the sequence moves from reachable in principle to deployable in fact.

\textbf{\rung{r1} (timing ceiling).} An oracle that controls only fetch timing, restricted to objects the predictor could name, recovers 95.8\% of the corridor on Wikipedia and 91.2\% on \texttt{cluster50}: these corridors are questions of \emph{when}, not \emph{which}. Given that the same oracle recovers 7.3\% on \texttt{cluster53}, its corridor is an object-choice problem, and that trace is withdrawn from the capture study on the pre-registered criterion. \emph{The capture study below therefore runs on two traces, not three}, a genuine restriction on the generality of its negative, reported here rather than left to the scope section. What supports the conclusion beyond trace count is that each rung fails for a separately measured reason and that the failures are mechanistically linked.

\textbf{\rung{r2} (arbitration).} Learning could help arbitration only if choosing fetches as a set were better than choosing them independently. Given identical value estimates and budget, an optimal joint chooser beats a per-candidate threshold rule by $+0.01$ and $-0.01$ points, far below the pre-registered 2-point requirement. The choice is separable; a threshold rule suffices, and reinforcement learning over this decision is unnecessary.

\textbf{\rung{r3} (the offered stream).} Perfect timing applied only to the candidates the baseline actually emits lands \emph{below} the baseline: $-7.01$ $[-7.68,-6.33]$ and $-10.82$ $[-11.70,-9.93]$. No scheduler can win by re-timing what is already offered, so the corridor is a problem of \emph{coverage}.

\textbf{\rung{r4} (coverage).} Widening the emission (fan-out 32, no confidence floor) and inserting just before use only those reuses the wide predictor could name in time clears the threshold: $+10.41$ $[8.74,12.08]$ and $+11.42$ $[9.95,12.91]$. The corridor is reachable in principle; breadth provides candidates and perfectly timed insertion converts them.

\textbf{\rung{r5} (emission-time fetch).} The natural causal version of \rung{r4} fetches each candidate when it is predicted and holds it until use. Granting this design a \emph{perfect} forecaster and \rung{r4}'s own selection rule, so that the only change is when the fetch happens, it still loses 10 to 15 points to \rung{r4} in every configuration, capturing at most 6\% of the corridor. Since a perfect forecaster fails when fetching at prediction time, no better forecaster can rescue the design; the missing ingredient is just-in-time insertion.

\textbf{\rung{r6} (causal placement).} \rung{r6} preserves \rung{r4}'s clairvoyant coverage and selection and replaces only the clock: the \texttt{oracle} arm fetches at the true next use and reproduces \rung{r4} ($+9.69$, $+11.69$, a passing construction check); the \texttt{model} arm fetches at a time predicted causally from recency, frequency, and confidence. The model meets the standard learnability gate (rank correlation $0.60$/$0.43$, never-reused AUC $0.92$/$0.81$) yet lands $-40.06$ $[-42.7,-37.4]$ and $-26.01$ $[-27.9,-24.1]$: precision collapses to $0.062$ and $0.054$. Wikipedia is the controlled case, since the model arm there issues essentially the baseline's own volume ($1.02$M against $0.98$M) and still loses 40 points; on \texttt{cluster50} volume grows to $1.49$M against $0.67$M. A pre-registered retrain on the wide stream raises rank correlation to $0.68$ without improving capture.

\begin{figure}[t]\centering
\figorbox{figures/ladder}{width=\columnwidth,height=9cm}
\caption{The Necessity Ladder as run (wiki / cluster50). Each rung prices one ingredient with a pre-registered oracle; green rungs certify the corridor reachable in principle, red rungs are negatives. The certificate reads trap: real, but reachable only with foresight.}
\label{fig:ladder}
\end{figure}

\begin{table}[t]
\centering\small
\setlength{\tabcolsep}{3pt}
\begin{tabular}{lllcc}
\toprule
selection & insertion & rung & wiki & c50 \\
\midrule
wide\orc & clairvoyant JIT\orc & \rung{r4} coverage & $+10.41$ & $+11.42$ \\
$k{=}1$ stream & clairvoyant JIT\orc & \rung{r3} stream & $-7.01$ & $-10.82$ \\
wide\orc & at emission & \rung{r5} emission & $+0.62$ & $-3.07$ \\
wide\orc & causal hazard & \rung{r6} placement & $-40.06$ & $-26.01$ \\
armed wide & causal event & \rung{r7} CGP & $-7.60$ & $-5.39$ \\
causal & causal & deployable & \multicolumn{2}{c}{$\le\min$ above} \\
\bottomrule
\end{tabular}
\caption{Capture decomposition (corridor versus baseline, points). \orc\ marks a clairvoyant ingredient. Every mixed row that removes one collapses; the all-causal rows sit far below the baseline. \rung{r5} is quoted in its \emph{most favourable} configuration, a perfect forecaster under a cumulative rather than instantaneous budget, so it is an upper bound on that design. \rung{r1} and \rung{r2} are reported in text.}
\label{tab:decomp}
\end{table}

\textbf{The mechanism: a break-even precision no coverage setting reaches.} Why does a model that ranks reuse well still fail? Rather than name a cause, the ladder prices one. A final arm holds coverage at \rung{r4}'s fan-out of 32, grants every genuinely reused candidate \rung{r4}'s own just-in-time insertion, and degrades nothing but precision, sweeping in a fraction of candidates that are never requested and issuing those at emission, where a policy believing in them would pay for them. The corridor is monotone in precision and crosses zero at a \textbf{break-even precision} of $p^{\star}=0.582$, between $+1.24$ $[+0.1,+2.4]$ at precision $0.629$ and $-3.04$ $[-3.9,-2.1]$ at $0.469$. A causal scheduler must therefore make hits out of 58\% of the bytes it consumes merely to match the tuned baseline at this coverage, and \rung{r6}'s causal arm converts $0.062$, short by a factor of $9.4$. Because perfect timing is granted outright here, the shortfall is a selection failure rather than a placement one; nor is it an eviction issue, since on the scored arm a prefetched object's hundred-request survival probability is $0.984$ against the baseline's $1.000$, and the implied demand-side hit rate is unchanged (supplementary document, \S C, which also records a mechanism claim these measurements withdraw). Narrowing does not recover it either, losing coverage faster than it repays precision: fan-out $1$ is $0.874$ precise and perfect timing on its own stream still loses ($\rung{r3}$, $-7.01$), while fan-outs $2$ through $8$ leave the causal arm flat near $-40$ as the \rung{r4} ceiling above it climbs from $-2.66$ to $+4.92$. No cut of this ranking is at once wide enough and precise enough to pay for its bandwidth: \textbf{a signal can be rank-learnable without being budget-feasible}. Pricing the gap in this manner also makes the negative independent of how well the timing model was built, which a failed arm alone cannot establish. Full sweeps are in the supplementary document, \S A--B.

\textbf{\rung{r7}: the constructive escape also fails.} A break-even precision suggests its own remedy, coverage without volume. Corroboration-Gated Prefetch (CGP) deposits wide-predictor candidates onto a byte-free watchlist and fires an actual fetch only when a tight predictor's signal triggers, choosing among watchlisted candidates by confidence within the baseline's exact budget. Coverage comes from the free watchlist, volume remains at the baseline's, and with perfect components the design reduces to \rung{r4}. It fails anyway: $-7.60$ $[-8.1,-7.1]$ and $-5.39$ $[-5.9,-4.9]$ against the tuned baseline. Since those figures compare a metered arm against an unmetered one, the capped-baseline control used in the LLM section is reported here too: against the baseline replayed under CGP's own token bucket the arms read $+0.98$ $[+0.93,+1.03]$ and $+1.01$ $[+0.94,+1.07]$, the bucket itself costing $-8.58$ and $-6.40$ points. Most of the raw shortfall is therefore the rate limit rather than the policy, and CGP still misses its pre-registered $2.5$-point gate under the like-for-like comparison. The reason is specific: raising the watchlist weight on \texttt{cluster50} moves the funded set from 43\% to 69\% watchlisted while precision stays at $0.268$, and on Wikipedia moves 41\% to 55\% at precision $0.577$ to $0.566$, the gate itself adding $+0.18$ and $-0.04$. Corroboration is statistically redundant with the predictor's own confidence, and its Wikipedia precision of $0.577$ approaches $p^{\star}$ only at roughly half \rung{r4}'s coverage, restating the frontier rather than escaping it. No causal policy in this design space secures coverage without paying precision below what the budget demands, a stronger negative than a simple failure because the policy was built against the exact bottleneck and the bottleneck absorbed it.

\textbf{What the ladder does and does not establish.} The bottom row of Table~\ref{tab:decomp}, that a deployable policy is bounded by the minimum above it, holds \emph{if} \rung{r3}--\rung{r7} bracket the deployable design space of selection and insertion. That is an inference from named instances plus a mechanism, not an impossibility proof: a design outside the decomposition, for instance one modifying the evictor to protect warmed objects, is not excluded by these measurements, though the substitutes result argues against it. What the rungs do establish is stronger than a single trained model's failure, because each removes one clairvoyant ingredient at a time and the collapse is attributable to the ingredient removed.

\textbf{What would have produced a build.} The Ladder was prepared to certify the opposite outcome: a build required one rung to hold under causal information. \rung{r4} shows the target exists and \rung{r6}'s oracle arm shows the scheduler design is sound, so a causal placement arm landing anywhere above the baseline would have read build, and a 2-point gap at \rung{r2} would have recommended a learned arbiter specifically. Neither occurred. The learned-index instantiation below exercises exactly this path and does return build, which is what distinguishes a procedure from a prior.

\section{The Procedure Across Domains}
A measurement procedure earns trust by serving different answers to different inputs. The identical arms and bars, without any logic modifications, were transferred over to two other settings, one manifesting as another negative answer and the other as a positive, and the two directions are subsequently inspected with trained models.

\textbf{LLM KV-cache serving: no corridor.} In the fixed one-hour traces of conversation and tool/agent traffic \citep{mooncake}, where objects are 512-token KV blocks, warming a block before next use is the prefetch analogue, and hit rate measures avoided prefill, the time-to-first-token lever. Two adaptations keep the port honest: warming a block is determined only at request boundaries, as per-access warming would fake-convert whole prefix chains at zero latency, and reachability is physical, as a block has to be computed before it can be warmed by any system. The answer is consistent under both comparators: against the tuned baseline the corridor is negative outright ($-0.56$, $-1.87$), and with the capped-baseline control providing a comparison more favourable to a corridor, the margin is $+2.98$ and $+2.14$ against the 8-point gate. A structural factor explains the verdict: prefix-chain reuse is near-deterministic (held-out coverage 0.728 versus 0.008 for popularity), which is why a tuned association table already sits on the physically reachable ceiling. The timing of reuse can be learned ($\rho{=}0.27$, AUC 0.81) on the conversation trace, yet there is no corridor for that signal to capture.

Learned indexes: a captured corridor. In a static learned index the oracle is the perfect index at one probe and the deployable arm is a two-stage recursive model (RMI) trained on the keys alone, completing with a last-mile search over its guaranteed error window [1, 16]. The only decisive difference from caching is reachability. A key is $N$ times its cumulative distribution and learnable from the keys themselves, so the oracle's advantage is based on data rather than the future and the procedure must be able to certify the result.
Certification requires sweeping both sides. Binary search is not the tuned non-learned baseline: a cache-conscious B-tree \citep{ccbtree} costs $\lceil\log_{B} N\rceil$ node accesses, and interpolation search \citep{interp} exploits precisely the position-is-CDF signal with no model at all, in $O(\log\log N)$ probes on smooth distributions. The baseline is therefore the best of $\{$binary, B-tree, interpolation$\}$ per dataset; symmetrically, the RMI's leaf count is swept, capped at $N/16$ so the model cannot degenerate into a position lookup table. That cap matches the B-tree's \emph{entry count}, not its bytes: at 24 B per leaf (slope, intercept, error bound) against 8 B per B-tree separator, $N/16$ leaves is $2.8\times$ the baseline's index memory. A byte-matched cap of $N/35$ is therefore reported alongside it, and the two survivors respond differently: \texttt{books} holds at 56.8\% while \texttt{meta\_rprn} falls from 50.4\% to 27.9\%, which is why only the former is claimed. Both cost units are reported, memory probes and distinct cache lines touched, because they rank these structures differently; the cache-line model treats the RMI root as resident and charges the leaf model and last-mile search, while the probe unit charges the RMI's last-mile search only and the B-tree every node access, a convention that favours the learned arm and is stated here rather than left implicit.

Table~\ref{tab:index} gives the verdicts: \textbf{build on one of five real key sets}, with a second marginal and sharp discrimination among the rest. A tuned \emph{RMI} beats the tuned family on \texttt{books} (67.5\% of the corridor by probes, 64.3\% by lines) and on the \texttt{meta\_rprn} identifiers (50.4\%/54.9\%), but the latter clears the bar by only 0.4 points and does so only under the entry-matched leaf cap: at byte parity it reads 27.9\% while \texttt{books} holds at 56.8\%, so only \texttt{books} is claimed. It is short on \texttt{osm\_cellids} and fails outright on \texttt{fb} and the Wikipedia identifiers. One measured quantity explains every verdict: how the \emph{maximum} error in each leaf, which sets the search window, changes with capacity. Where the bulk is within a cache line the corridor is captured; where capacity helps but runs out it is approached and missed (\texttt{osm\_cellids} drops to $7.05$ probes, never reaching the B-tree's $6.00$); where the model is uniformly ill-fitted capacity is inert, \texttt{fb} and \texttt{wiki\_2019t} holding 100\% of keys beyond sixty-four lines across a 125-fold leaf increase, a structural failure of piecewise-linear approximation which no budget can fix. The two cost units also disagree in both directions across the suite, so a single-metric evaluation would have returned a wrong verdict twice; a correctness invariant, every key inside its declared window, holds in every run. Per-key-set distributions are in the supplementary document, \S E.

\textbf{The sweep rule reverses verdicts in both directions.} This suite is where the discipline is most visible, because applying one rule symmetrically both destroyed and restored conclusions. Scoring the learned index against binary search alone, the field's customary comparison, is mechanism M1 exactly; adding the two competitors the rule requires removed up to 290 points of apparent capture and reversed three of five verdicts. Applying the same rule to the learned side then reversed two of them back, because sweeping the baseline while pinning the RMI at one leaf count is the identical error inverted: at fixed capacity \texttt{books} reads 41.7\% and \texttt{meta\_rprn} 10.4\%, both below bar, and once the model is swept as well they read 67.5\% and 50.4\%, both above it. A procedure that only ever removed positives would be a rejection machine; one that removes and restores under a single rule is a measurement.


\begin{table}[t]
\centering\footnotesize
\setlength{\tabcolsep}{2pt}
\begin{tabular}{llrrrl}
\toprule
key set & keys & base & probes & lines & verdict \\
\midrule
\texttt{books} & 2M & btree & 67.5 & 64.3 & \textbf{build} \\
\texttt{meta\_rprn} & 1.0M & interp & 50.4 & 54.9 & marginal$^{\dagger\ddagger}$ \\
\texttt{osm\_cellids} & 2M & btree & $-$21.0 & 30.2 & reach., unc. \\
\texttt{fb} & 2M & btree & $-$140.0 & $-$66.7 & reach., unc. \\
\texttt{wiki\_2019t} & 0.80M & btree & $-$267.3 & $-$111.5 & reach., unc. \\
\midrule
synth.\ uniform & 2M & interp & 47.2 & 50.2 & unit-dep. \\
synth.\ lognormal & 2M & btree & $-$5.5 & 43.9 & reach., unc. \\
synth.\ clustered & 2M & btree & $-$73.6 & $-$11.4 & reach., unc. \\
\bottomrule
\end{tabular}
\caption{Learned indexes: percentage of the lookup-cost corridor (tuned baseline to perfect index) captured by a deployable RMI under two cost units, both sides swept. \emph{build} clears the pre-registered 50\% bar on both units. Reachability holds by construction here (position is $N\cdot\mathrm{CDF}$), so the negative rows are not \emph{no corridor} but a fourth category the caching settings never exhibit: \emph{reachable but uncaptured by this model class}. $^\dagger$Clears by 0.4 points. $^\ddagger$Fails at byte parity. Key-set provenance and both sensitivity analyses are in the supplementary document, \S E; the calibration of the positive rate is discussed in \S G.}
\label{tab:index}
\end{table}

\textbf{A substrate-competence control.} A procedure characterised only where it says build has an unmeasured false-negative rate, and that is the expensive error: advising against a component that would have worked is a mistake nobody observes. Two trained cache-management policies were therefore run, each frozen after training on its trace's first half and strictly causal at decision time: a gradient-boosted reuse-distance evictor like LRB \citep{lrb} and a gradient-boosted admission classifier like Baleen \citep{baleen} with a tuned bypass threshold. Because neither spends speculative bandwidth, neither is subject to the break-even precision above; the trained prefetch-side test is \rung{r6} itself, which is, and which lands $-40.06$/$-26.01$. They instead check whether a learned component of an entirely \emph{different} capability incidentally captures what the procedure declared unreachable, and whether the substrate admits learned wins at all. Each is scored against the \textbf{eviction-management corridor} (Belady over S3-FIFO with no prefetch anywhere), distinct from the prefetch corridor of Table~\ref{tab:learnable}: $+11.55$, $+12.50$, and $+12.75$ on \texttt{wiki\_2019t}, \texttt{cluster50}, and \texttt{cluster53}. Against those the trained models reach essentially none, the evictor landing $1.83$ points below the baseline on Wikipedia and the admission classifier capturing $-0.0\%$ and $0.0\%$ on the two key-value traces. The latter two are a tuning outcome rather than a capture measurement and are reported as such: the bypass threshold is tuned on the training half, and tuning selected the degenerate setting, admit everything, collapsing the policy onto S3-FIFO exactly. The failures are informative because the same evictor is competent where the substrate allows, beating S3-FIFO by $+3.45$, $+3.13$, and $+0.99$ points on three traces whose corridors fall below the bar, every interval excluding zero; they are therefore properties of the corridors, not of the learners.

\textbf{The recipe.} The instantiations above follow one procedure, stated so it can be run on a new domain. Fix the deployable baseline by sweeping the strongest non-learned family, tuned against the corridor's interest; meter every oracle arm to that baseline's measured resource budget and require Pareto dominance; limit each oracle to information a deployable model could hold, and charge the restriction inside the replay rather than by subtraction; pre-register thresholds and invariants before the first real run, reporting every arm that the escalation kills; and, if a corridor survives, price each capability a learner could add with a separate oracle, certifying \emph{build} only if an all-causal combination clears the bar.


\noindent Step 3 is what headroom studies miss, and Step 5 is what separates certification from measurement: a corridor only a clairvoyant arm can reach is a trap, and the recipe says so before a model is trained.

\section{Scope and Threats to Validity}
``Learnable'' is defined with respect to a frozen, history-based predictor. An online-updating or content-aware predictor could reach cold objects and would raise the surviving corridors, so the measurement verdicts are conservative in that direction; the capture verdicts are not symmetric, since \rung{r4} already grants the widest emission the family supports and \rung{r6}/\rung{r7} fall short on precision at that width, which a larger vocabulary does not remedy. The two remaining limits on the capture negative, its two-trace scope and the spanning assumption of Table~\ref{tab:decomp}, are made explicit where they arise; the measurement verdicts rest on all twelve traces. A 1\% cache over 2M-request prefixes is the primary operating point, and cache-size and byte-weighted analyses point in the same direction. Twitter traces are one-in-ten samples, LLM traces span one hour of one provider, and intervals cover within-trace variance only. Every error in the baseline inflates corridors, so the negatives are conservative under a closer look.

\section{Conclusion}
Measured against a family-tuned baseline at matched bandwidth, most celebrated prefetch headroom in caching is phantom: 14--87\% of gross corridors evaporate across twelve production traces. What survives is real but, priced by pre-registered oracle decompositions, reachable only with foresight, at a quantified price, a break-even precision of $0.582$ against $0.062$ achieved, and in a general form, rank-learnable is not budget-feasible. The procedure's value is that it discriminates: across four settings it returns \emph{no corridor} (block storage, LLM KV serving), a \emph{trap} (CDN and key-value prefetching), and a \emph{build} (learned indexes), each decided by replay before training, with a captured corridor and a competence control bounding it on both sides. Reachability proves necessary but not sufficient, since the green light is withheld wherever a tuned classical structure already sits on the reachable ceiling. The negatives also point somewhere, and now with a target attached: the binding constraint is precision at width, not learning capacity, so a break-even precision converts a negative result into a specification a later system can be measured against. The procedure and every replay arm are provided as anonymized supplementary material, so the next celebrated headroom number can be checked in an afternoon.

\bibliography{references}

\end{document}
